\scp{\tsc{\ili{Ambon}-Seram}}{C-05}
Some \tsc{\ili{Ambon}-Seram} languages are spoken in areas where two
kinds of marsupials are found: cuscus (\it{Phalangeridae}) throughout,
and bandicoots (\it{Peroryctidae}) in the \ili{east}.

\sscp{Cuscus \it{(Phalangeridae)}}{C-05-02}
Terms for ‘cuscus’ (or ambiguous ‘marsupial’) in \tsc{\ili{Ambon}-Seram}
languages are given in the following table.
Note that while we faithfully reproduce
\citeauthor{Collins1983}'s (\citeyear{Collins1983})
vague gloss of `marsupial' which he applies to his *ma(n)sər and its reflexes,
wherever more detailed information is
available these terms universally designate `cuscus'.

Similarly, \citet{Stresemann1927} uses
\it{Baummarder} `tree marten' and \it{Beutelmarder} `pouch marten'
interchangeably for Proto-\ili{Ambon} **maḏer.
We translate these terms `cuscus', as that is the designation
whenever independent evidence is available.
Thus, for instance \citet[41]{Stresemann1927} has
``Beutelmarder S.A. [Sub-\ili{Ambon}] *maḏel-ĕ > \ili{Paulohi} \it{makel-e}''
while \citet[130]{Stresemann1918} gives the definition of \ili{Paulohi}
as ``cuscus \it{Phalanger orientalis} and \it{Ph. maculatus}''.

Several terms for kinds of cuscus are from PCM **sibaw
‘cuscus pouch, female cuscus’ which is based on from \ili{Watui} \it{sibaʔwe} ‘marsupial sp.’,
Naka{\Q}ela \it{sihaw-e} ‘marsupial sp.’,
\ili{Hulung} \it{sipa} ‘marsupial sp.’,
\ili{Yalahatan} \it{mare siha} ‘female (gray) cuscus’,
South \ili{Nuaulu} \it{mara siha {\tl} siha} ‘cuscus (kind,
female)’, \ili{Buru} \it{sifa-n} ‘marsupial pouch,
cuscus pouch’.

\newpage
\noindent{\wg\st{2pt}\begin{xltabular}{\linewidth}
{lllll>{\raggedright\arraybackslash}X}
\lsptoprule
%Language	&	ISO	&	Term	&	Gloss	&	Etymon	&	Source	\\	\midrule
%\endfirsthead
Language	&	ISO	&	Term	&	Gloss	&	Etymon	&	Source	\\	\midrule
\endhead
East \ili{Alune}	&	alp	&	marel,	&	cuscus,	&	**mantəR,	&	\cite[196]{Ellen1993},	\\	\hhline{~}
	&		&	sibakw-e	&	female cuscus	&	**sibaw 	&	\cite[298]{Niggemeyer1951b}	\\
West \ili{Alune}	&	alp	&	maʧel	&	cuscus	&	**mantəR	&	\cite[305]{vanEkris1864}	\\
Naka{\Q}ela	&	nae	&	sihaw-e	&	marsupial sp.	&	**sibaw	&	\cite[49]{Collins1983}	\\
\ili{Hulung}	&	huk	&	sipa	&	marsupial	&	**sibaw	&	\cite[49]{Collins1983}	\\
\ili{Kayeli}	&	kzl	&	maren	&	cuscus	&	**mantəR	&	\cite[40]{Stresemann1927}	\\
S. \ili{Wemale}	&	tlw	&	madel-e	&	marsupial	&	**mantəR	&	\cite[55]{Collins1983}	\\
N. \ili{Wemale}	&	weo	&	marel-e	&	marsupial	&	**mantəR	&	\cite[55]{Collins1983}	\\
\ili{Batumerah}	&		&	marel-a	&	cuscus	&	**mantəR	&	\cite[254]{Ludeking1868} \\ %\cite[86]{Collins1983}	\\
\ili{Manipa}	&	mqp	&	marel-e	&	marsupial	&	**mantəR	&	\cite[68]{Collins1983}	\\
\ili{Piru}	&	ppr	&	marel	&	cuscus	&	**mantəR	&	\cite[305]{vanEkris1864}	\\
\ili{Larike}	&	alo	&	márid-u	&	cuscus	&	**mantəR	&	\cite[77]{LaidigLaidig1991}	\\
\ili{Wakasihu}	&	alo	&	maril-u	&	marsupial	&	**mantəR	&	\cite[68]{Collins1983}	\\
\ili{Alang}	&	alo	&	maril-u	&	marsupial	&	**mantəR	&	\cite[84]{Collins1983}	\\
\ili{Asilulu}	&	asl	&	marel	&	cuscus	&	**mantəR	&	\cite[70]{Collins2003b}	\\
\ili{Paulohi}	&	plh	&	makel-e	&	cuscus	&	**mantəR	&	\cite[130]{Stresemann1918}, \cite[196]{Ellen1993}	\\
\ili{Kaibobo}	&	kzb	&	makel	&	cuscus	&	**mantəR	&	\cite[305]{vanEkris1864}	\\
\ili{Hitu}	&	htu	&	makel	&	cuscus	&	**mantəR	&	\cite[68]{Collins1983}, \cite[196]{Ellen1993}	\\
\ili{Kamarian}	&	kzx	&	maker	&	cuscus	&	**mantəR	&	\cite[305]{vanEkris1864}	\\
\ili{Haruku}	&	hrk	&	makel	&	cuscus	&	**mantəR	&	\cite[305]{vanEkris1864}	\\
\ili{Kailolo}	&	hrk	&	makel-e	&	marsupial	&	**mantəR	&	\cite[68]{Collins1983}	\\
\ili{Saparua}	&	spr	&	makel	&	cuscus	&	**mantəR	&	\cite[305]{vanEkris1864}	\\
\ili{Latu}	&	ltu	&	makell-o	&	marsupial	&	**mantəR	&	\cite[68]{Collins1983}	\\
\ili{Nusalaut}	&	nul	&	makel	&	cuscus	&	**mantəR	&	\cite[305]{vanEkris1864}	\\
\ili{Amahai}	&	amq	&	maker-o	&	cuscus	&	**mantəR	&	\cite[43]{Stresemann1918}	\\
\ili{Yalahatan}\footnotemark 	&	jal	&	marer-e,	&	cuscus,	&	**mantəR,	&	\cite{BeckleyND}	\\	\hhline{~}
	&		&	mare siha	&	female cuscus	&	**sibaw	&		\\
\end{xltabular}}
\footnotetext{Fuller cuscus-related information for \ili{Yalahatan} from Beckley
		(n.d.): \it{marer-e} ‘cuscus (generic)’; \it{mare maliate} ‘spotted
		cuscus’, \it{mare malote} ‘male gray cuscus’,
		\it{mare siha} ‘female gray cuscus’,
		\it{mare elaʔe} ‘large cuscus sp.’ (cf. \it{elaʔe} ‘large’).
		Loss of a root-final C{\#} in the presence of attributive modifiers
		is common in many languages across Wallacean subgroups.}

\noindent{\wg\st{2pt}\begin{xltabular}{\linewidth}
{lllll>{\raggedright\arraybackslash}X}
\lsptoprule
Language	&	ISO	&	Term	&	Gloss	&	Etymon	&	Source	\\	\midrule
\endhead
S. \ili{Nuaulu}	&	nxl	&	maran-e,	&	cuscus,	&	**mantəR,	&	\cite[196]{Ellen1993}	\\	\hhline{~}
	&		&	siha	&	female cuscus	&	**sibaw	&	\cite[65, 102]{Bolton-Grimes1990,MatokeBolton2005}\\
\ili{Manusela}	&	wha	&	ihisi	&	cuscus	&	\it{misc.}	&	\cite{MacDonald2015}	\\
\ili{Liana-Seti}	&	ste	&	mair-a	&	cuscus	&	**mantəR	&	\cite{McCollum-Grimes1990}	\\
\ili{Benggoi}	&	bgy	&	med-am	&	cuscus	&	**mantəR	&	Nicholas Williams (pers. comm. July 2024)	\\
\ili{Hoti}	&	hti	&	mala	&	cuscus	&	**mantəR	&	\cite[101]{Stresemann1927}	\\
\lspbottomrule
\end{xltabular}}

\sscp{Bandicoots \it{(Peroryctidae)}}{C-05-01}
Terms for bandicoots are currently only known from two languages,
South \ili{Nuaulu} and \ili{Manusela}, given below.
Regarding the South \ili{Nuaulu} Terms Roy Ellen (pers. comm. July 2024) doubts
that \it{imanona} is ‘bandicoot’. He states “I collected lots of non-marsupial rats that were
confidently identified by people as \it{imanona}. However, descriptions of \it{tui-tui} come
close to what we know of \it{R. prattorum}, though of course \it{tui-tui} may be a synonym of
other terms that I never discovered.''
\citet{MatokeBolton2005} gloss \it{tui-tui} as `civet'
(Indonesian \it{musang, tenggalung}, \ili{Ambon} \ili{Malay} \it{tinggalong}).
It is one of three terms, alongside \it{kuha} and \it{makwene} with this gloss.
\mbox{}\\

\noindent{\wg\st{2pt}\begin{tabularx}
{\linewidth}{>{\hspace{-\tabcolsep}}l<{\hspace{-\tabcolsep}}
>{\hspace{-\tabcolsep}\enspace}ll>{\raggedright\arraybackslash}Xl}
\lsptoprule
%Language 	&	ISO	&	Term	&	Gloss	&	Source	\\	\midrule
%\endfirsthead
Language 	&	ISO	&	Term	&	Gloss	&	Source	\\	\midrule
%\endhead
S. \ili{Nuaulu}	&	nxl	&	imanona	&	marsupial (long nose, smaller than cuscus);					&	\cite{Bolton-Grimes1990};	\\
					&			&					&	big rat;																										&	\cite[43]{MatokeBolton2005}	\\
					&			& tui-tui	& Seram long-nosed bandicoot, \it{Rhynchomeles prattorum} ?;	& Roy Ellen pers. comm. July 2024; \\
					&			& 				& civet 																										& \cite[114]{MatokeBolton2005} \\
\ili{Manusela}	&	wha	&	mape-a	&	Seram long-nosed bandicoot, \it{Rhynchomeles prattorum}		&	\cite{MacDonald2015}	\\
\lspbottomrule
\end{tabularx}\mbox{}\\}